\documentclass[12pt,letterpaper]{article}
\usepackage{amsmath}
\usepackage{amsthm}
\usepackage{amsfonts}
\usepackage{amssymb}
\usepackage{amscd}
\usepackage{enumerate}
\usepackage{fancyhdr}
\usepackage{mathrsfs}
\usepackage{bbm}
\usepackage{framed}
\usepackage{mdframed}
\usepackage{cancel}
\usepackage{mathtools}
\usepackage{verbatim}
\usepackage{hyperref}
\usepackage[letterpaper,voffset=-.5in,bmargin=3cm,footskip=1cm]{geometry}
\setlength{\parindent}{0.0in}
\setlength{\parskip}{0.1in}
\allowdisplaybreaks
\headheight 15pt
\headsep 10pt
\newcommand\N{\mathbb N}
\newcommand\Z{\mathbb Z}
\newcommand\R{\mathbb R}
\newcommand\Q{\mathbb Q}
\newcommand\lcm{\operatorname{lcm}}
\newcommand\setbuilder[2]{\ensuremath{\left\{#1\;\middle|\;#2\right\}}}
\newcommand\E{\operatorname{E}}
\newcommand\V{\operatorname{V}}
\newcommand\Pow{\ensuremath{\operatorname{\mathcal{P}}}}

\DeclarePairedDelimiter\ceil{\lceil}{\rceil}
\DeclarePairedDelimiter\floor{\lfloor}{\rfloor}
\newcommand\hint[1]{\textbf{Hint}: #1}
\newcommand\note[1]{\textbf{Note}: #1}
\newcounter{enum22i}
\newenvironment{22enumerate}{\begin{enumerate}[a.]\itemsep0em\setcounter{enumi}{\value{enum22i}}}{\setcounter{enum22i}{\value{enumi}}\end{enumerate}}
\newenvironment{22itemize}{\begin{itemize}\itemsep0em}{\end{itemize}}
\fancypagestyle{firstpagestyle} {
  \renewcommand{\headrulewidth}{0pt}
  \lhead{\textbf{CSCI 0220}}
  \chead{\textbf{Discrete Structures and Probability}}
  \rhead{M. Littman}
}

\fancypagestyle{fancyplain} {
  \renewcommand{\headrulewidth}{0pt}
  \lhead{\textbf{CSCI 0220}}
  \chead{Midterm}
  \rhead{\textit{ }}
}
\pagestyle{fancyplain}
\usepackage{tikz}

\begin{document}
  \thispagestyle{firstpagestyle}
  \begin{center}
    {\huge \textbf{CSCI0220 Midterm}}

    {\large Due: \textit{June 29th at 11:59pm EDT}}
  \end{center}
  
 {\large \textbf{Please read in full:}}
 
You are not allowed to use the internet outside of the CSCI0220 website. Doing so is a breach of the \textbf{Collaboration Policy} and \textbf{Brown's Academic Code}.

You can refer to lecture slides, lecture captures, your own notes, the course textbook, the website resources, and Campuswire. You must complete the entire exam by yourself, without discussion of any kind with anyone else.

\textbf{By submitting this exam, you are confirming that you have understood these policies, and that you have followed them.}

As always, please do not include any identifying information about yourself in the handin, including your Banner ID. Be sure to fully explain your reasoning and show all work for full credit.


\subsection*{Problem 1}
{
\nopagebreak 

{\setcounter{enum22i}{0}
\begin{22enumerate}

    \item Prove using the set-element method that for finite sets $X$ and $Y$ with the same universal set $U$

    \[X \cap Y = (X \cup Y)-(\overline{X} \cup \overline{Y}) .\]
    
\item \begin{enumerate}[i.]
    \item Consider an arbitrary binary operation $\star$ on a finite set $X$. A \textit{left identity element} for $\star$ on the set $X$ is any $e \in X$ such that $e \star x = x$ for all $x \in X$. For example, $1$ is the left identity element for the multiplication operator `$\times$' on the integers $\Z$ because $1 \times x = x$ for all $x\in\Z$. 
    
\bigskip

Find the left identity elements of the set difference operator `$-$' on the set $\Pow(S)$, where $S$ is an arbitrary non-empty, finite set. If there exists one, \textbf{prove} why the element you chose is an identity element,\textbf{ and prove} why no other element would be. If there does not exist one, \textbf{prove} that no element can be.

\item By the same logic, a \textit{right identity element} for $\star$ on the set $X$ is any $e\in X$ such that $x\star e=x$ for all $x\in X$. For example, $1$ is the right identity element for the division operator `$/$' on the integers $\Z$ because $x/1=x$ for all $x\in\Z$.

\bigskip

Find the right identity elements of the set difference operator `$-$' on the same set $\Pow(S)$. If there exists one, \textbf{prove} why the element you chose is an identity element,\textbf{ and prove} why no other element would be. If there does not exist one, \textbf{prove} that no element can be.
\end{enumerate}
\end{22enumerate}
	
    
}
}


\subsection*{Problem 2}
{
{\setcounter{enum22i}{0}

    Let $B$ be the set of all possible propositional formulas. Note that $B$ is a set of infinite cardinality. Define the relation $R$ on $B$ such that, for $a,b \in B$, $a\;R\;b$ if and only if $a$ and $b$ always return the same output given the same input.

	\begin{22enumerate}	

	\item
	Let $x$ be a proposition with Boolean inputs $p$, $q$ and $r$ and the following truth table:
	\begin{center}
	\begin{tabular}{c | c | c || c}
	$p$ & $q$ & $r$ & $x$\\ \hline
	1 & 1 & 1 & 1\\
	1 & 1 & 0 & 0\\
	1 & 0 & 1 & 0\\
	1 & 0 & 0 & 0\\
	0 & 1 & 1 & 1\\
	0 & 1 & 0 & 0\\
	0 & 0 & 1 & 1\\
	0 & 0 & 0 & 0\\
	\end{tabular}
	\end{center}

	Write a propositional formula corresponding to $x$ in conjunctive normal form with a maximum of two clauses.\\
	\textbf{Note: }Explain how you arrived at your answer.

	\item
	Prove that $R$ is an equivalence relation.

	
	\item Create a bijection between the equivalence classes of three-input, one-output propositional formulas and binary strings of length eight.\\
	\textbf{Note: }You need not prove your answer. You must just explain the reasoning behind your bijection.
	
	\item How many equivalence classes of three-input, one-output
	propositional formulas are there?\\
	\textbf{Note: }You must thoroughly explain your answer, but are not required to prove it.
	
	\item 
	The equivalence class of a propositional formula $x$ is $[x]_R = \{a \in B \mid (a,x) \in R\}$. Prove that $|[x]_R|$ is infinite for any $x \in B$.
	\end{22enumerate}
}
}



\subsection*{Problem 3}
{
\nopagebreak 

{\setcounter{enum22i}{0}
Let the relation $R_{15} = \{(a,b) \;|\; \exists i \in \Z \text{ such that } (a-b) = 15i\}$. Prove by induction that, for any finite set $S$ with more than one element, $(\left|\Pow(\Pow(S))\right|,1)\in R_{15}$.
}
}


\end{document}